\section{Discussion}
\label{sec:discussion}

\subsection{Cross-View Aggregation of Program Reasoning}
The forward and backward tasks are two views of one object: the same program, the same gold value,
the same grader, asked from opposite ends. That is what makes their disagreement informative. An arm
that improves the forward view and degrades the backward one on the same programs has not become
better at reading those programs; it has become better at producing the answer it was trained to
produce. Reading each view alone hides this --- every adaptation in Table~\ref{tab:main} looks like
progress forward --- and only aggregating across views exposes it.

The aggregate we propose is simple and strict: an adaptation should be at or above the untuned model
in \emph{every} view, on \emph{every} model, before it is called an improvement in understanding
rather than in task fit. Under that criterion exactly one arm in this study passes on all eight
models, the sign-consensus merge (Section~\ref{sec:results}), and it does so without being the best
arm in any single view. Two properties of this criterion are worth stating. It is conservative in the
right direction: it can only remove arms from the winner's list, never add one. And it costs nothing
to apply beyond evaluating in more than one direction, which any behaviour-prediction benchmark can
do, since the pairing is a property of the items rather than of the model.

\subsection{Choosing an Adaptation Strategy}

No arm dominates, and the table is more useful than a ranking. If the deployment distribution is known
and fixed, the anchored objective and the mixture are the strongest forward arms and the three to six
points are real. If it is not --- if an obfuscator may be updated, transforms may be composed in
unseen ways, or the model may be asked something other than what it was tuned on --- the merge is the
only arm in this study that is never significantly below the untuned model in any block on any model,
and it costs one adapter and no inference-time machinery. The mixture matches or beats it on the unseen
family but needs eight adapters and a gate at inference, and its learned gate is worth nothing over a
uniform one on the held-out family.

\subsection{The Limits of Training Breadth}
The intuition behind pooled training is that a model shown more transforms learns something more
general. The results say the opposite in both out-of-distribution regimes. On the held-out family,
pooled training is below every arm composed from the same data on all eight models and below an
adapter that never saw an obfuscated program on six. Backwards, it is below the untuned model on four
of the six models where it is readable. Neither is what an invariance would predict; both are what
memorising a transform's surface predicts, since the surface is exactly what an unseen obfuscator is
free to change.

The point is not that breadth is useless --- on seen compositions it is within a few points of the
best arm --- but that its gains are confined to the distribution that produced them, and that the
same data, learned as separate specialists and composed afterwards, does not have that confinement.
The difference between the two is not the data or the capacity; it is whether the conditions were
learned in one set of weights or in six. Where the deployment distribution is known and fixed that
distinction may not matter. Where an obfuscator can be updated, it is the whole question.
